% How the LaTeX preview works: a short Beamer talk, to open in Sagebrush's LaTeX editor (press Present for the slides).
% Written by Claude Opus 5.5, an AI model made by Anthropic, for the Sagebrush project (SageMath, Inc.).
% Same license as the rest of articles/. Numbers: bench/latex, November 4, 2026 (see 01-how-it-works.md).
\documentclass{beamer}
\usetheme{Madrid}

\title[The LaTeX preview]{How the \LaTeX{} preview works}
\author[Sagebrush]{Claude Opus 5.5, for the Sagebrush project}
\date{November 4, 2026}

\begin{document}

\begin{frame}
  \titlepage
\end{frame}

\section{Two parts}

\begin{frame}{An expander, then a renderer}
  \begin{itemize}
    \item<1-> \textbf{The expander} reads the source as \TeX{} does: macros,
      groups, registers, conditionals, with bounds on work.
    \item<2-> \textbf{The renderer} turns its events into HTML, with
      mathematics by KaTeX: $\displaystyle\sum_{k=1}^{n} k = \frac{n(n+1)}{2}$.
    \item<3-> Both run in a Web Worker, so typing never waits.
  \end{itemize}
\end{frame}

\section{Testing}

\begin{frame}{Real \TeX{} as the oracle}
  Compare the words of \TeX{}'s PDF, $P$, with the preview's, $Q$:
  \[
    \operatorname{Dice}(P, Q) = \frac{2\,|P \cap Q|}{|P| + |Q|}
  \]
  \pause
  A low value is the signal; the words missing most say why.
\end{frame}

\section{Results}

\begin{frame}{A typed character to the preview}
  \begin{center}
  \begin{tabular}{lrr}
    \hline
    document & \texttt{latexmk} & preview \\
    \hline
    this kind of talk & 3.88 s & 32 ms \\
    pgfplots manual & 21 s & 258 ms \\
    arXiv paper & 801 ms & 100 ms \\
    \hline
  \end{tabular}
  \end{center}
  \medskip
  Medians of 5 runs, November 4, 2026.
\end{frame}

\begin{frame}{What would have saved time}
  \begin{enumerate}
    \item Real \TeX{} as the oracle from the first day.
    \item Measure what the reader sees.
    \item Ask \TeX{} instead of guessing.
    \item Report what you can't show, from the start.
  \end{enumerate}
\end{frame}

\end{document}
